% Target: rmp-workbench-iii-mpo-injective-white
% Formula: A virtual MPO corner pulls through a white PEPS tensor unchanged.
% Ink: Canonical public kernel PEPS star with before-and-after corner paths.
% Boundary: Five PEPS legs and two MPO string ends remain open on each side.
% Source: author source ImagesReview Section III/ImagesReview Section III.tex
%   lines 501-520; archival output MPOinjnog.pdf
% Capabilities: kernel, pulling-through, open-string
\[
  % Author panel (lines 501-518): the virtual MPO corner pulled through the
  % PEPS tensor.  The g-string crosses the west and back legs with an MPO
  % tensor at each crossing; pulled through, it crosses the east and front
  % legs instead, and the tensor A is unchanged.
  \tenkzkernel
  \tndeclare{species}{lattice}{hue=source:black}
  \tndeclare{species}{g}{hue=source:blue}
  \begin{tenkzeq}[size=m]
    \begin{tenkz}[rows={wire,wire,wire}, cols=3, bonds=none]
      % Four virtual legs end at A; its typed physical port stays source-short.
      \tn[at=(2,2), name=A,
        ports={0:virtual,45:virtual,90:physical,180:virtual,225:virtual}]{}
      \tnwire[kind=index, species=lattice, name=vxw]{(2,1)}{A.180}
      \tnwire[kind=index, species=lattice, name=vxe]{A.0}{(2,3)}
      \tnwire[kind=index, species=lattice, name=vdsw]{(3,1)}{A.225}
      \tnwire[kind=index, species=lattice, name=vdne]{A.45}{(1,3)}
      % the MPO corner crosses the west and back legs, one dot per crossing
      \tn[at=midway (1,1) and (2,2), skin=none, name=jg1]{}
      \tn[at=midway (2,2) and (1,3), skin=none, name=jg2]{}
      \tnwire[kind=string, species=g, name=g1,
              cross={over at crossing of g1 and vxw}]
        {midway (2,1) and (3,1)}{jg1}
      \tnwire[kind=string, species=g, name=g2,
              cross={over at crossing of g2 and vdne,
                     over at crossing of g1 and g2}]
        {jg1}{jg2}
      \tnwire[kind=string, species=g, name=g3,
              cross={over at crossing of g3 and vdne}]
        {jg2}{midway (1,3) and (2,3)}
      \tn[species=g, at=crossing of g1 and vxw]{}
      \tn[species=g, at=crossing of g2 and vdne]{}
      \tnmark[form=label, label pos=nw]{on g1 0.8195}{$g$}
      \tnmark[form=label, label pos=se]{A}{$A$}
    \end{tenkz}
    =
    \begin{tenkz}[rows={wire,wire,wire}, cols=3, bonds=none]
      \tn[at=(2,2), name=A,
        ports={0:virtual,45:virtual,90:physical,180:virtual,225:virtual}]{}
      \tnwire[kind=index, species=lattice, name=vxw]{(2,1)}{A.180}
      \tnwire[kind=index, species=lattice, name=vxe]{A.0}{(2,3)}
      \tnwire[kind=index, species=lattice, name=vdsw]{(3,1)}{A.225}
      \tnwire[kind=index, species=lattice, name=vdne]{A.45}{(1,3)}
      % the same corner, pulled through to the east and front legs
      \tn[at=midway (2,2) and (3,3), skin=none, name=jg1]{}
      \tn[at=midway (2,2) and (3,1), skin=none, name=jg2]{}
      \tnwire[kind=string, species=g, name=g1,
              cross={over at crossing of g1 and vxe}]
        {midway (1,3) and (2,3)}{jg1}
      \tnwire[kind=string, species=g, name=g2,
              cross={over at crossing of g2 and vdsw,
                     over at crossing of g1 and g2}]
        {jg1}{jg2}
      \tnwire[kind=string, species=g, name=g3,
              cross={over at crossing of g3 and vdsw}]
        {jg2}{midway (2,1) and (3,1)}
      \tn[species=g, at=crossing of g1 and vxe]{}
      \tn[species=g, at=crossing of g2 and vdsw]{}
      \tnmark[form=label, label pos=se]{on g2 0.5837}{$g$}
      \tnmark[form=label, label pos=nw]{A}{$A$}
    \end{tenkz}
  \end{tenkzeq}.
\]
